% ============================================================================
%  SECCIÓN 3.5: REQUERIMIENTOS PRELIMINARES
% ============================================================================

\section{Requerimientos preliminares}

Los requerimientos describen lo que el sistema debe hacer (requerimientos
funcionales) y las restricciones o cualidades que debe cumplir
(requerimientos no funcionales) \citep{sommerville2011, kendall2011}. A
continuaci\'on se presentan los requerimientos preliminares de la aplicaci\'on
m\'ovil, que servir\'an de base para el desarrollo del MVP.

\subsection{Requerimientos funcionales (RF)}

\begin{table}[H]
  \centering
  \caption{Requerimientos funcionales preliminares}
  \contenidoConFuente{%
    \begin{tablaAPA}
      \begin{tabular}{@{}p{2.0cm}p{3.4cm}p{10.4cm}@{}}
        \toprule
        \textbf{ID} & \textbf{Nombre} & \textbf{Descripci\'on} \\
        \midrule
        RF-01 & Registro de usuarios & El sistema permitir\'a a los usuarios
        registrarse indicando nombres, correo electr\'onico y contrase\~na, y
        seleccionar su perfil (asistente, gestor o administrador). \\
        \addlinespace
        RF-02 & Inicio de sesi\'on & El sistema permitir\'a a los usuarios iniciar
        sesi\'on con su correo y contrase\~na, y cerrar sesi\'on de forma segura. \\
        \addlinespace
        RF-03 & Cat\'alogo de escenarios & El sistema mostrar\'a el listado de
        escenarios deportivos registrados con su informaci\'on b\'asica. \\
        \addlinespace
        RF-04 & B\'usqueda y filtros & El sistema permitir\'a buscar escenarios por
        nombre, deporte, ciudad y tipo de escenario. \\
        \addlinespace
        RF-05 & Mapa interactivo & El sistema mostrar\'a la ubicaci\'on de los
        escenarios en un mapa georreferenciado y permitir\'a consultar los m\'as
        cercanos a la ubicaci\'on del usuario. \\
        \addlinespace
        RF-06 & Detalle del escenario & El sistema mostrar\'a la ficha de cada
        escenario con sus caracter\'isticas, disciplinas, horarios, servicios y
        eventos. \\
        \addlinespace
        RF-07 & Gesti\'on de favoritos & El sistema permitir\'a al usuario guardar
        escenarios como favoritos y consultar la lista de sus favoritos. \\
        \addlinespace
        RF-08 & Gesti\'on de contenido & El sistema permitir\'a a los gestores
        registrar, editar y desactivar la informaci\'on de sus escenarios. \\
        \addlinespace
        RF-09 & Notificaciones & El sistema enviar\'a notificaciones sobre eventos
        y novedades de los escenarios guardados como favoritos. \\
        \addlinespace
        RF-10 & Administraci\'on & El sistema permitir\'a al administrador gestionar
        usuarios y validar la informaci\'on publicada por los gestores. \\
        \bottomrule
      \end{tabular}
    \end{tablaAPA}%
  }{Elaboraci\'on propia en base a la revisi\'on documental.}
\end{table}

\subsection{Requerimientos no funcionales (RNF)}

\begin{table}[H]
  \centering
  \caption{Requerimientos no funcionales preliminares}
  \contenidoConFuente{%
    \begin{tablaAPA}
      \begin{tabular}{@{}p{2.0cm}p{3.4cm}p{10.4cm}@{}}
        \toprule
        \textbf{ID} & \textbf{Atributo} & \textbf{Descripci\'on} \\
        \midrule
        RNF-01 & Usabilidad & La interfaz deber\'a ser intuitiva y de f\'acil uso
        para usuarios sin experiencia previa, siguiendo principios de dise\~no
        centrado en el usuario \citep{nielsen1994}. \\
        \addlinespace
        RNF-02 & Rendimiento & El sistema deber\'a responder a las consultas en un
        tiempo m\'aximo de 3 segundos en condiciones normales y el mapa deber\'a
        mostrarse de forma fluida en dispositivos m\'oviles. \\
        \addlinespace
        RNF-03 & Seguridad & La autenticaci\'on y el manejo de datos personales
        deber\'an cumplir las pr\'acticas de seguridad est\'andar, protegiendo la
        informaci\'on de los usuarios. \\
        \addlinespace
        RNF-04 & Compatibilidad & La aplicaci\'on m\'ovil deber\'a funcionar en los
        sistemas operativos iOS y Android. \\
        \addlinespace
        RNF-05 & Accesibilidad & La interfaz deber\'a considerar pautas de
        accesibilidad que permitan su uso por personas con discapacidad,
        incluyendo contraste adecuado y textos alternativos. \\
        \addlinespace
        RNF-06 & Disponibilidad & El sistema deber\'a estar disponible al menos el
        99\,\% del tiempo durante el periodo de evaluaci\'on. \\
        \addlinespace
        RNF-07 & Portabilidad & El c\'odigo deber\'a estar organizado por m\'odulos
        para facilitar su mantenimiento y la incorporaci\'on de nuevas
        funcionalidades. \\
        \addlinespace
        RNF-08 & Calidad de software & El sistema deber\'a satisfacer los atributos
        de calidad establecidos en el modelo ISO/IEC 25010, como funcionalidad,
        eficiencia y mantenibilidad \citep{iso25010}. \\
        \bottomrule
      \end{tabular}
    \end{tablaAPA}%
  }{Elaboraci\'on propia en base a la revisi\'on documental.}
\end{table}
